\documentclass{article}

\usepackage{microtype}
\usepackage{graphicx}
\usepackage{booktabs}
\usepackage{hyperref}
\newcommand{\theHalgorithm}{\arabic{algorithm}}

\usepackage{icml2025}

\usepackage{amsmath}
\usepackage{amssymb}
\usepackage{amsthm}
\usepackage{mathtools}
\usepackage{multirow}
\usepackage{array}
\usepackage{enumitem}
\usepackage{subcaption}

\icmltitlerunning{Compositional Feature Interpretability for Diffusion Language Models}

\newcommand{\LR}{\mathrm{LR}}
\newcommand{\bfx}{\mathbf{x}}
\newcommand{\bfc}{\mathbf{c}}
\newcommand{\bfb}{\mathbf{b}}
\newcommand{\bfh}{\mathbf{h}}
\newcommand{\bfbeta}{\boldsymbol{\beta}}
\newcommand{\bfgamma}{\boldsymbol{\gamma}}
\newcommand{\softplus}{\mathrm{softplus}}

\begin{document}

\twocolumn[
\icmltitle{Compositional Feature Interpretability for Diffusion Language Models}

\begin{icmlauthorlist}
\icmlauthor{Shijian Xu}{unibas}
\end{icmlauthorlist}

\icmlaffiliation{unibas}{Department of Mathematics and Computer Science, University of Basel, Basel, Switzerland}
\icmlcorrespondingauthor{Shijian Xu}{shijian.xxu@gmail.com}

\icmlkeywords{mechanistic interpretability, sparse autoencoders, diffusion language models, feature compositionality, nonlinear interactions}

\vskip 0.3in
]

\printAffiliationsAndNotice{}

\begin{abstract}
Diffusion language models (dLLMs) denoise text through bidirectional masked attention, raising an open question: how do SAE-discovered features compose with each other as denoising progresses?
We apply NIMO---a nonlinear interaction model with a closed-form linear decomposition---to sparse autoencoder (SAE) features extracted from LLaDA-8B-Base~\citep{llada2025} across six transformer layers and five denoising timesteps.
To handle the high-dimensional, ultra-sparse feature space ($d_\mathrm{SAE} = 16{,}384$, top-$K = 50$), we develop \emph{SparseSAENIMO}, a compressed context-encoding adaptation of NIMO that scales to SAE dimensions without the $O(d^2)$ memory cost of the original formulation.
The Linearity Ratio ($\LR$)---the fraction of model output attributable to linear feature effects---decreases monotonically as the timestep decreases (more context becomes available) and as layer depth increases.
At the deepest measured layer (30) and smallest timestep ($t = 0.10$), $\LR = 0.094$, indicating that $90.6\%$ of the model's log-probability signal at that point is mediated by nonlinear cross-feature interactions.
In contrast, a causal autoregressive baseline (GPT-2 Small~\citep{radford2019language} with Joseph Bloom SAEs~\citep{bloom2024sae}) shows $\LR \ge 0.30$ across all layers and prefix lengths, confirming that bidirectional attention generates qualitatively stronger feature interactions than unidirectional attention.
These results constitute the first quantitative characterisation of compositional feature dynamics in a large-scale diffusion language model.
\end{abstract}

%% ─────────────────────────────────────────────────────────────────
\section{Introduction}
%% ─────────────────────────────────────────────────────────────────

Sparse autoencoders (SAEs) have emerged as a powerful tool for mechanistic interpretability, decomposing dense neural representations into sparse, approximately monosemantic features~\citep{cunningham2023sparse,templeton2024scaling}.
However, identifying individual features is only the first step: understanding how features \emph{compose}---how they interact to produce model outputs---is equally important and less studied.

Diffusion language models (dLLMs)~\citep{llada2025,sahoo2024simple,lou2024discrete} present a particularly interesting setting for this question.
Unlike autoregressive models, which compute representations through unidirectional attention, dLLMs operate via bidirectional masked attention across all tokens simultaneously at each denoising step.
This architecture creates a rich environment for cross-feature interactions: each feature is computed in the context of all other token features in the sequence.
Moreover, the denoising process introduces a controllable variable---the noise level $t$---that varies from fully masked ($t = 1.0$, no context) to nearly clean ($t \to 0$, full context).
This provides a natural axis along which to study how context availability shapes feature compositionality.

NIMO (Nonlinear Interaction Model Order)~\citep{xu2026nimo} provides a principled framework for measuring this compositionality.
NIMO decomposes a model's prediction into
\begin{equation}
    \hat{y} = \beta_0 + \sum_j x_j\,\beta_j\!\left(1 + g_u(\bfx_{-j},\,\mathbf{E}_j)\right),
    \label{eq:nimo}
\end{equation}
where $\beta_j$ captures the global linear effect of feature $j$ and $g_u(\bfx_{-j}, \mathbf{E}_j)$ captures how the prediction from feature $j$ is modulated by the remaining context $\bfx_{-j}$.
The Linearity Ratio
\begin{equation}
    \LR = \frac{\|\bfbeta\|_1}{\|\bfbeta\|_1 + \mathbb{E}_n\!\left[\|G_n^{\mathrm{active}}\|_1\right]}
    \label{eq:lr}
\end{equation}
quantifies the fraction of the model's output attributable to linear (context-independent) effects. $\LR = 1$ means purely linear; $\LR = 0$ means purely nonlinear.

We make the following contributions:
\begin{enumerate}[leftmargin=*]
    \item \textbf{SparseSAENIMO.} A NIMO variant that scales to SAE feature spaces ($d_\mathrm{SAE} = 16{,}384$) by replacing the $O(d^2)$ context tensor with a compressed binary positional encoding of dimension $2 \times 15 = 30$, making the correction network input dimension independent of $d_\mathrm{SAE}$.

    \item \textbf{Phase transition in $\LR(t)$.} LR decreases monotonically as $t$ decreases, across all six LLaDA-8B-Base layers tested. The transition is strongest in deep layers: layer~30 goes from $\LR = 0.413$ at $t = 1.0$ (degenerate) to $\LR = 0.094$ at $t = 0.10$.

    \item \textbf{Depth--compositionality gradient.} LR decreases with layer depth: shallow layer~1 ($\LR \in [0.79, 0.88]$) is substantially more linear than deep layer~30 ($\LR \in [0.09, 0.41]$).

    \item \textbf{AR vs.\ dLLM contrast.} GPT-2 Small shows flat $\LR \ge 0.30$ across all prefix lengths in early/mid layers, confirming that the strong nonlinearity in dLLM is architecturally specific to bidirectional masked attention.

    \item \textbf{Feature categorisation.} Across all layers, roughly 450--490 features are ``always-on'' (stable across timesteps), while 560--1,200 are ``early-dominant'' and 670--1,200 are ``late-dominant'', revealing a structured turnover in the feature vocabulary.
\end{enumerate}

%% ─────────────────────────────────────────────────────────────────
\section{Background}
%% ─────────────────────────────────────────────────────────────────

\subsection{Diffusion Language Models}

LLaDA-8B-Base~\citep{llada2025} is an 8B-parameter masked diffusion language model based on a bidirectional transformer (32~layers, $d_\mathrm{model} = 4{,}096$).
Given a clean token sequence, LLaDA corrupts it by independently masking each token with probability $t$, then trains a bidirectional transformer to predict the original tokens from the masked sequence.
At generation time, LLaDA iteratively unmasks tokens over $T$ denoising steps.
At noise level $t$, the model conditions on all visible tokens simultaneously, creating dense bidirectional attention patterns that differ fundamentally from autoregressive models~\citep{sahoo2024simple,austin2021structured}.

\subsection{Sparse Autoencoders for LLMs}

DLM-Scope~\citep{dlmscope2024} provides pre-trained TopK SAEs for LLaDA-8B-Base at layers $\{1, 6, 11, 16, 26, 30\}$.
Each SAE encodes a residual stream activation $\bfh \in \mathbb{R}^{4096}$ as
\begin{equation}
    \mathbf{f} = \mathrm{TopK}(\mathbf{W}_{\mathrm{enc}}\,\bfh + \mathbf{b}_{\mathrm{enc}},\;K=50) \;\in\; \mathbb{R}^{16384},
\end{equation}
retaining only the $K = 50$ largest activations~\citep{gao2024scaling}. The resulting feature vectors are ultra-sparse: only 50 out of 16,384 dimensions are nonzero per sample (sparsity ratio 99.7\%).

\subsection{NIMO}

NIMO~\citep{xu2026nimo} is a nonlinear interaction model that combines a linear backbone with a context-conditioned correction:
\begin{equation}
    \hat{y}_k = \gamma_{0k} + \sum_{j \in \mathcal{V}} \tilde{x}_{jk}\,\gamma_{jk},
    \qquad
    \tilde{x}_{jk} = x_j\,c_{jk}\!\left(1 + g_u(\bfx_{-j}, \mathbf{E}_j)\right),
\end{equation}
where $c_{jk} = \softplus(C_{jk}) > 0$ is a learnable positive scale, $g_u$ is a shared MLP, and $\gamma$ is computed in closed form by ridge regression on the effective features $\tilde{\mathbf{x}}$.
The interpretable linear coefficient is $\beta_{jk} = c_{jk}\,\gamma_{jk}$, and the nonlinear correction $G_{njk} = g_u(\bfx_{-j,n}, \mathbf{E}_j)$ measures how much the context modulates feature $j$'s contribution for sample $n$.

%% ─────────────────────────────────────────────────────────────────
\section{SparseSAENIMO}
%% ─────────────────────────────────────────────────────────────────

\subsection{The Scalability Problem}

Direct application of NIMO to SAE features is infeasible: computing the leave-one-out context $\bfx_{-j}$ for all $V = 2{,}048$ vocabulary features per sample requires building an $[N, V, d_\mathrm{SAE}]$ tensor, which at $N = 27{,}000$, $V = 2{,}048$, $d_\mathrm{SAE} = 16{,}384$ would require approximately 900\,GB per batch.

\subsection{Compressed Context Encoding}

We replace $\bfx_{-j} \in \mathbb{R}^{d_\mathrm{SAE}}$ with a binary-position-encoded summary:
\begin{align}
    \mathrm{context\_sum}_i &= \sum_{j'\ \mathrm{active}} f_{j',i} \cdot \bfb_{j'} \;\in\; \mathbb{R}^{n_\mathrm{bits}}, \\
    \bfc_{-j,i} &= \mathrm{context\_sum}_i - f_{j,i} \cdot \bfb_j \;\in\; \mathbb{R}^{n_\mathrm{bits}},
\end{align}
where $\bfb_j \in \{-0.5, +0.5\}^{n_\mathrm{bits}}$ is the centered binary representation of index $j$, and $n_\mathrm{bits} = \lfloor \log_2 d_\mathrm{SAE} \rfloor + 1 = 15$ for $d_\mathrm{SAE} = 16{,}384$.
The correction network receives $[\bfc_{-j} \mid \bfb_j] \in \mathbb{R}^{30}$ as input.

\paragraph{Key property.}
$\bfc_{-j} = \mathbf{0}$ when all other features are zero, so $g_u(\mathbf{0}, \mathbf{E}_j) = 0$ by construction (after bias centering), preserving NIMO's linear-when-isolated property.

\paragraph{Memory footprint.}
The computation now requires only an $[N, V, 30]$ tensor, reducing memory by a factor of $d_\mathrm{SAE} / (2 \times n_\mathrm{bits}) \approx 546\times$.

\subsection{Correction Network Architecture}

The correction network $g_u$ is a three-layer MLP with skip connections and sinusoidal activations:
\begin{align}
    z_1 &= \tanh\!\bigl(0.3 \cdot \mathrm{fc}_1([\bfc_{-j} \mid \bfb_j])\bigr), \\
    z_2 &= \sin\!\bigl(2\pi \cdot \mathrm{fc}_2([z_1 \mid \bfb_j])\bigr), \\
    G   &= \mathrm{fc}_3([z_2 \mid z_1 \mid \bfb_j]).
\end{align}
The feature identity $\bfb_j$ is re-injected at layers 2 and 3 to prevent it from being washed out by the context aggregation. Input noise ($\sigma = 0.1$) is added to $z_1$ during training for regularisation.

\subsection{Vocabulary Selection}

We scan the training set and select the $V_\mathrm{active} \le V_\mathrm{max} = 2{,}048$ most frequently activated SAE features. NIMO's L1 penalty via the $C$ regulariser further shrinks most $\beta$ coefficients toward zero.

At $t = 1.0$ (fully masked input), only 50 distinct features are ever active (those of the mask token embedding). This near-degenerate regime is treated as a boundary condition.

\subsection{Corrected Linearity Ratio}
\label{sec:lr}

The naive LR formula sums $|G[n, v, k]|$ for all $V$ vocabulary features. However, at any given sample only $K \approx 50$ features are nonzero; summing over inactive features overcounts the nonlinear term by a factor of $V/K \approx 41\times$, artificially deflating LR. We correct this:
\begin{equation}
    \LR = \frac{\|\bfbeta\|_1}{\|\bfbeta\|_1 + \mathbb{E}_n\!\left[\displaystyle\sum_{j:\,f_{nj} \neq 0} |G_{nj}|\right]}.
    \label{eq:lr_corrected}
\end{equation}
This \emph{active-feature LR} measures the ratio of linear to total signal only over features that actually contribute to each prediction.

\subsection{Training}

Each (layer, timestep) NIMO is trained independently with:
Adam optimiser ($\mathrm{lr} = 3\times10^{-4}$, cosine annealing over 60 epochs), batch size 512, regularisation $\lambda = \mu = 0.5$, full-dataset $\gamma$ re-solve after every epoch, and a 90/10 train/validation split with random seed~42.

%% ─────────────────────────────────────────────────────────────────
\section{Experimental Setup}
%% ─────────────────────────────────────────────────────────────────

\subsection{Dataset}

We sample 30,000 sequences from the Pile~\citep{gao2020pile} validation split (512 tokens each) and run LLaDA-8B-Base forward at noise levels $t \in \{0.10, 0.25, 0.50, 0.75, 1.00\}$.
At each $t$, we randomly mask each token independently with probability $t$, then record:
\begin{itemize}
    \item \textbf{SAE features}: residual stream activations at layers $\{1, 6, 11, 16, 26, 30\}$, encoded through DLM-Scope TopK SAEs $\to$ shape $[N, d_\mathrm{SAE}]$ with at most 50 nonzeros per row.
    \item \textbf{Log-probability target}: the model's mean log-probability over masked token positions.
\end{itemize}
This yields 30 feature files (6 layers $\times$ 5 timesteps), each with $N = 30{,}000$ samples.

\subsection{AR Baseline}

For comparison, we extract features from GPT-2 Small~\citep{radford2019language} (12 layers, $d_\mathrm{model} = 768$) using Joseph Bloom's SAEs~\citep{bloom2024sae} ($d_\mathrm{SAE} = 24{,}576$) at layers $\{0, 2, 4, 6, 8, 10\}$.
We vary the prefix fraction $p \in \{0.10, 0.25, 0.50, 0.75, 1.00\}$ (the fraction of tokens visible), giving 19,998 samples per (layer, prefix) pair.
The AR prefix fraction $p$ is the natural analogue of the dLLM noise level $t$: both control how much context is available.

\subsection{Evaluation Metrics}

\begin{itemize}
    \item \textbf{Linearity Ratio (LR)}: primary metric, Eq.~\eqref{eq:lr_corrected}.
    \item \textbf{Validation $R^2$}: fraction of log-prob variance explained by NIMO.
    \item \textbf{Linear probe $R^2$}: ridge regression baseline (same vocabulary subset).
    \item \textbf{Beta stability}: feature categorisation across timesteps (always-on, early-dominant, late-dominant).
    \item \textbf{Interaction density}: fraction of feature pairs with significant $G$ modulation.
\end{itemize}

%% ─────────────────────────────────────────────────────────────────
\section{Results}
%% ─────────────────────────────────────────────────────────────────

\subsection{Phase Transition in LR($t$)}

\begin{table}[t]
\centering
\caption{Linearity Ratio $\LR(t)$ for LLaDA-8B-Base. Rows marked $*$ are degenerate ($t = 1.0$, only 50 mask-token features active) and are excluded from trend analysis.}
\label{tab:lr_dlm}
\vskip 0.1in
\begin{small}
\begin{sc}
\begin{tabular}{r cccc c}
\toprule
\textbf{Layer} & $t{=}0.10$ & $t{=}0.25$ & $t{=}0.50$ & $t{=}0.75$ & $\Delta$ \\
\midrule
1  & 0.790 & 0.814 & 0.855 & 0.878 & $+0.088$ \\
6  & 0.624 & 0.659 & 0.719 & 0.775 & $+0.151$ \\
11 & 0.527 & 0.559 & 0.623 & 0.680 & $+0.153$ \\
16 & 0.434 & 0.462 & 0.531 & 0.584 & $+0.150$ \\
26 & 0.199 & 0.224 & 0.259 & 0.303 & $+0.104$ \\
30 & 0.094 & 0.108 & 0.124 & 0.135 & $+0.041$ \\
\bottomrule
\end{tabular}
\end{sc}
\end{small}
\end{table}

Table~\ref{tab:lr_dlm} shows that LR increases monotonically as $t$ increases for all six layers ($t = 0.10 \to 0.75$).
This is consistent with the phase transition hypothesis: as more tokens are masked (larger $t$), the model has less bidirectional context, reducing cross-feature interactions and making the log-probability signal more linearly predictable from individual features.

LR also decreases with layer depth at any fixed $t$. At $t = 0.10$, LR drops from 0.790 (layer~1) to 0.094 (layer~30)---a ratio of $8.4\times$. This suggests that deeper layers operate in a substantially more nonlinear compositional regime, consistent with the view that early layers perform local feature detection while later layers perform nonlinear feature combination~\citep{elhage2021mathematical}.

\subsection{Feature Beta Stability Across Timesteps}

We categorise each feature by its $|\beta(t)|$ trajectory across timesteps within a layer:
\begin{itemize}
    \item \textbf{Always-on}: stable contribution magnitude across timesteps (relative std $< 0.3$).
    \item \textbf{Early-dominant}: stronger at small $t$ (more context $\to$ more active).
    \item \textbf{Late-dominant}: stronger at large $t$ (less context $\to$ more active).
\end{itemize}

\begin{table}[t]
\centering
\caption{Feature category counts per layer ($V = 2{,}048$ vocabulary; threshold $|\beta| > 10^{-3}$).}
\label{tab:beta_stability}
\vskip 0.1in
\begin{small}
\begin{sc}
\begin{tabular}{r rrrr}
\toprule
\textbf{Layer} & \textbf{AO} & \textbf{ED} & \textbf{LD} & \textbf{Total} \\
\midrule
1  & 483 & 1,038 & 1,203 & 2,724 \\
6  & 487 &   840 & 1,116 & 2,443 \\
11 & 456 &   935 & 1,148 & 2,539 \\
16 & 452 &   931 & 1,156 & 2,539 \\
26 & 461 &   841 & 1,102 & 2,404 \\
30 & 388 &   563 &   671 & 1,622 \\
\bottomrule
\end{tabular}
\end{sc}
\end{small}
\vskip -0.05in
{\small AO: always-on; ED: early-dominant; LD: late-dominant.}
\end{table}

The majority of active features change their importance across timesteps (Table~\ref{tab:beta_stability}): only 19--22\% are always-on in layers 1--26 (24\% in layer~30).
Late-dominant features (those more influential when context is sparse) slightly outnumber early-dominant features at most layers, suggesting that many features serve as fallback predictors when context is unavailable.
Layer~30 has notably fewer active features (1,622 vs.\ $\sim$2,500 in other layers), consistent with a more compressed vocabulary at the deepest layer.

\subsection{Feature Rankings and Causal Relevance}

For each (layer, timestep) pair we compute three feature ranking criteria: (a)~$|\beta_j|$ (NIMO global importance), (b)~activation frequency (proportion of samples where feature $j$ is active), and (c)~absolute Pearson correlation with the log-probability target.

The top-10 $|\beta|$ features shift substantially across timesteps within a layer. In layer~26, the leading feature at $t = 0.10$ (SAE index~6185, $|\beta| = 0.053$) is replaced by a different feature (index~861, $|\beta| = 0.071$) at $t = 0.75$, indicating that the most linearly predictive feature changes as context availability varies. Full rankings are saved in \texttt{results/causal\_ranking.pt} for downstream activation-patching experiments~\citep{meng2022locating}.

\subsection{Feature Interaction Sparsity}

For each layer we fit a sparse linear interaction model on the $G$ corrections:
\begin{equation}
    G_{ij} \approx \sum_{j'} \alpha_{j'j}\,f_{j'},
\end{equation}
keeping the top-200 edges by $|\alpha_{j'j}|$.
In all six layers, 200 significant interactions correspond to $\le 0.005\%$ of all $V^2 = 2{,}048^2 = 4{,}194{,}304$ possible feature pairs.
This extreme sparsity is consistent with the mechanistic interpretability literature suggesting that SAE features interact through small circuits rather than all-to-all dependencies~\citep{elhage2021mathematical,conmy2023towards}.

\subsection{dLLM vs.\ AR-LLM Compositionality}

\begin{table}[t]
\centering
\caption{Linearity Ratio $\LR$ for GPT-2 Small (AR) across prefix fractions. ``Flat?'' indicates std across prefix fractions $< 0.05$.}
\label{tab:lr_ar}
\vskip 0.1in
\begin{small}
\begin{sc}
\begin{tabular}{r ccccc r}
\toprule
\textbf{L} & $p{=}0.10$ & $p{=}0.25$ & $p{=}0.50$ & $p{=}0.75$ & $p{=}1.00$ & \textbf{Flat?} \\
\midrule
0  & 0.975 & 0.975 & 0.973 & 0.971 & 0.970 & Yes \\
2  & 0.610 & 0.591 & 0.588 & 0.598 & 0.561 & Yes \\
4  & 0.564 & 0.523 & 0.522 & 0.532 & 0.538 & Yes \\
6  & 0.538 & 0.484 & 0.468 & 0.480 & 0.499 & Yes \\
8  & 0.505 & 0.439 & 0.363 & 0.377 & 0.389 & No  \\
10 & 0.457 & 0.371 & 0.301 & 0.303 & 0.302 & No  \\
\bottomrule
\end{tabular}
\end{sc}
\end{small}
\end{table}

AR layers 0--6 show flat LR across prefix fractions ($\mathrm{std} < 0.05$; Table~\ref{tab:lr_ar}), confirming that causal attention creates stable, context-independent feature effects.
Layers 8 and 10 show modest variation ($\mathrm{std} \approx 0.052$--$0.061$), suggesting that deep layers in causal models also develop some prefix-length sensitivity.

Critically, the minimum AR LR is 0.301 (layer~10, $t = 0.50$), while the minimum dLLM LR is 0.094 (layer~30, $t = 0.10$)---a $3.2\times$ gap. At equivalent relative depth (${\sim}83\%$), dLLM layer~26 achieves $\LR = 0.199$ vs.\ AR layer~10's $\LR = 0.301$, a $1.5\times$ gap. This difference is architecturally attributable to bidirectional attention: dLLM features see the entire sequence simultaneously at each denoising step, creating richer cross-feature modulation.

\subsection{NIMO vs.\ Linear Baseline}

\begin{table}[t]
\centering
\caption{NIMO validation $R^2$ vs.\ linear probe $R^2$ on log-probability prediction (selected layers/timesteps, excluding $t = 1.0$ degenerate).}
\label{tab:r2}
\vskip 0.1in
\begin{small}
\begin{sc}
\begin{tabular}{rr cc r}
\toprule
\textbf{L} & $t$ & \textbf{Lin.} & \textbf{NIMO} & \textbf{Impr.} \\
\midrule
1  & 0.10 & 0.0467 & 0.0488 & $+4.5\%$ \\
1  & 0.75 & 0.1431 & 0.1480 & $+3.4\%$ \\
16 & 0.10 & 0.1776 & 0.1789 & $+0.7\%$ \\
16 & 0.75 & 0.2362 & 0.2405 & $+1.8\%$ \\
26 & 0.10 & 0.2993 & 0.3042 & $+1.6\%$ \\
26 & 0.75 & 0.3005 & 0.3137 & $+4.4\%$ \\
30 & 0.10 & 0.3016 & 0.3034 & $+0.6\%$ \\
30 & 0.75 & 0.2590 & 0.2661 & $+2.7\%$ \\
\bottomrule
\end{tabular}
\end{sc}
\end{small}
\end{table}

NIMO consistently outperforms the linear probe (Table~\ref{tab:r2}), confirming that the nonlinear correction $g_u$ captures genuine signal beyond linear feature effects.
The improvement is modest (0.6--4.5\%) but consistent across all 28 non-degenerate (layer, $t$) pairs, validating the model's nonlinear component.
Absolute $R^2$ values are moderate (0.05--0.36), reflecting that log-probability is a complex aggregate target partially explained by any single layer's features.

%% ─────────────────────────────────────────────────────────────────
\section{Discussion}
%% ─────────────────────────────────────────────────────────────────

\paragraph{Phase transition interpretation.}
The monotonic LR$(t)$ curve has a natural mechanistic interpretation. At large $t$ (many masked tokens), each token's representation is computed in a sparse context---most attention keys come from mask tokens with near-constant representations. The model therefore cannot leverage cross-token feature interactions, so log-probability is approximately linear in individual feature activations. As $t$ decreases, real tokens fill the context, bidirectional attention can integrate diverse features, and log-probability becomes a genuinely nonlinear function of the joint feature activation pattern.

The depth gradient amplifies this effect: deeper layers have processed the bidirectional context through more attention heads and MLP layers, accumulating more complex feature interactions. By layer~30, even at $t = 0.75$ (only 25\% context masked), $\LR = 0.135$, showing that deep features in LLaDA are almost never linearly predictive in isolation.

\paragraph{$t = 1.0$ degeneracy.}
At $t = 1.0$, every input token is masked, so the SAE sees only the mask token embedding (a single fixed vector, repeated across positions). All 30,000 samples produce near-identical feature activations, collapsing the vocabulary to exactly $K = 50$ active features. NIMO training diverges and LR values at $t = 1.0$ are unreliable (e.g., $\LR = 0.997$ for layer~11---a numerical artifact of near-constant features). This is not a bug but an informative boundary: at $t = 1.0$, LLaDA cannot use context because there is none.

\paragraph{AR deeper layer non-flatness.}
Layers 8 and 10 of GPT-2 Small show mild prefix-fraction dependence ($\mathrm{std} \approx 0.06$). This suggests that even in autoregressive models, the deepest layers develop some context sensitivity---likely because longer prefixes provide more input tokens for deep attention heads to integrate. However, the effect is 3--10$\times$ smaller than the corresponding dLLM effect at similar relative depth.

\paragraph{Limitations.}
\begin{itemize}
    \item \textbf{Single dLLM architecture.} Results are specific to LLaDA-8B-Base with DLM-Scope SAEs. Whether other dLLM architectures (e.g., MDLM~\citep{sahoo2024simple}, Plaid~\citep{gulrajani2024likelihood}) show the same phase transition remains an open question.
    \item \textbf{Log-probability as target.} We predict token-level log-probability as a proxy for semantic information content. Other targets (e.g., next-timestep feature activation, generation quality metrics) may reveal different LR profiles.
    \item \textbf{Compressed context.} The binary positional encoding summarises feature co-activation patterns but cannot represent actual context feature values. A richer encoding (e.g., learned feature embeddings) may capture stronger interactions.
    \item \textbf{Causal steering not yet executed.} The feature rankings in \texttt{results/causal\_ranking.pt} enable activation patching experiments, but these experiments are not yet implemented.
\end{itemize}

%% ─────────────────────────────────────────────────────────────────
\section{Conclusion}
%% ─────────────────────────────────────────────────────────────────

We have applied NIMO to SAE features from LLaDA-8B-Base to characterise how feature compositionality evolves along the denoising trajectory and across layer depth.
The central finding is a robust phase transition: the Linearity Ratio decreases as the noise level decreases and as layer depth increases, reaching $\LR = 0.094$ at layer~30 / $t = 0.10$.
An autoregressive baseline (GPT-2 Small) shows $\LR \ge 0.30$ under comparable conditions, confirming that bidirectional masked attention generates qualitatively stronger nonlinear feature interactions.

These results suggest that mechanistic interpretability methods developed for autoregressive models---which largely assume approximately additive feature effects~\citep{elhage2022toy}---may underestimate interaction complexity when applied to diffusion language models. Methods that explicitly account for nonlinear feature interactions, such as NIMO, are better suited to characterising the computational structure of dLLMs.

%% ─────────────────────────────────────────────────────────────────
\bibliographystyle{icml2025}
\bibliography{paper}

%% ─────────────────────────────────────────────────────────────────
\appendix

\section{Implementation Details}
\label{app:impl}

\paragraph{SparseSAENIMO hyperparameters.}
$d_\mathrm{SAE} = 16{,}384$; $n_\mathrm{bits} = 15$; $V_\mathrm{max} = 2{,}048$; $\mathrm{hidden\_dim} = 64$; $\lambda = \mu = 0.5$; epochs $= 60$; batch size $= 512$; $\mathrm{lr} = 3\times10^{-4}$; seed $= 42$.

\paragraph{Hardware.}
Two NVIDIA L40S GPUs (46\,GB VRAM each).
Feature extraction: ${\sim}1.1$~hours (30,000 sequences $\times$ 5 timesteps $\times$ 6 layers on a single GPU).
NIMO fitting: ${\sim}4$~hours total (30 models distributed across both GPUs).
Analysis: ${<}5$~minutes.

\paragraph{Software.}
PyTorch~2.x; \texttt{AwesomeInterpretability/llada-mask-topk-sae} for DLM-Scope SAEs; \texttt{jbloom/GPT2-Small-SAEs-Reformatted} for AR SAEs.

\section{Full LR Tables}
\label{app:full_lr}

\begin{table}[h]
\centering
\caption{Full LR table for LLaDA-8B-Base, including degenerate $t=1.0$ rows.}
\label{tab:full_lr_dlm}
\vskip 0.1in
\begin{small}
\begin{sc}
\begin{tabular}{r ccccc r}
\toprule
\textbf{L} & $t{=}0.10$ & $t{=}0.25$ & $t{=}0.50$ & $t{=}0.75$ & $t{=}1.00^*$ & $R^2$ \\
\midrule
1  & 0.790 & 0.814 & 0.855 & 0.878 & 0.523 & 0.049 \\
6  & 0.624 & 0.659 & 0.719 & 0.775 & 0.969 & 0.070 \\
11 & 0.527 & 0.559 & 0.623 & 0.680 & 0.996 & 0.155 \\
16 & 0.434 & 0.462 & 0.531 & 0.584 & 0.952 & 0.179 \\
26 & 0.199 & 0.224 & 0.259 & 0.303 & 0.081 & 0.304 \\
30 & 0.094 & 0.108 & 0.124 & 0.135 & 0.414 & 0.303 \\
\bottomrule
\end{tabular}
\end{sc}
\end{small}
\vskip -0.05in
{\small $^*$Degenerate: only 50 mask-token features active; LR unreliable.}
\end{table}

\begin{table}[h]
\centering
\caption{Full LR table for GPT-2 Small (AR baseline).}
\label{tab:full_lr_ar}
\vskip 0.1in
\begin{small}
\begin{sc}
\begin{tabular}{r ccccc}
\toprule
\textbf{L} & $p{=}0.10$ & $p{=}0.25$ & $p{=}0.50$ & $p{=}0.75$ & $p{=}1.00$ \\
\midrule
0  & 0.975 & 0.975 & 0.973 & 0.971 & 0.970 \\
2  & 0.610 & 0.591 & 0.588 & 0.598 & 0.561 \\
4  & 0.564 & 0.523 & 0.522 & 0.532 & 0.538 \\
6  & 0.538 & 0.484 & 0.468 & 0.480 & 0.499 \\
8  & 0.505 & 0.439 & 0.363 & 0.377 & 0.389 \\
10 & 0.457 & 0.371 & 0.301 & 0.303 & 0.302 \\
\bottomrule
\end{tabular}
\end{sc}
\end{small}
\end{table}

\end{document}
